\chapter{Mark-Katalog}
\label{app:marks}

Dieser Anhang enthält alle verfügbaren Prägungen.

\section{Abstammungen}\label{sec:abstammungen}
\noindent\textbf{Asgoran}\par
\begin{itemize}[leftmargin=*,nosep]
\item Vorteil: Erhalte +10 auf Proben mit Geschick + Athletik oder Geist + Mythen, die mit Wasser, Schiffen oder Küsten zu tun haben.
\item Verwundbarkeit: In engem Inlandsgelände (enge Gassen, tiefe Höhlen) erhältst du -10 auf Bewegungsproben bei Verfolgung oder Flucht.
\end{itemize}
\noindent\textbf{Doldagor}\par
\begin{itemize}[leftmargin=*,nosep]
\item Vorteil: Einmal pro Szene gilt vertikale Bewegung für dich als trivial, und du erhältst +10 auf Proben zur Positionierung in der Höhe.
\item Verwundbarkeit: Nach einem misslungenen Manöver mit hoher Beweglichkeit erleidest du eine zusätzliche geringe Konsequenz (Aufprall, Bloßstellung oder Desorientierung).
\end{itemize}
\noindent\textbf{O'Grut}\par
\begin{itemize}[leftmargin=*,nosep]
\item Vorteil: Erhalte +10 auf Täuschungsproben und soziale Irreführung, wenn Identität oder Absicht verschleiert sind.
\item Verwundbarkeit: Wirst du öffentlich direkt herausgefordert, erhältst du -10 auf die erste unmittelbare soziale Abwehrprobe, bis du klar Stellung beziehst.
\end{itemize}
\noindent\textbf{Flit}\par
\begin{itemize}[leftmargin=*,nosep]
\item Vorteil: Erhalte +10 auf Proben zum Ausweichen oder Umpositionieren in offenem Raum und ignoriere einmal pro Szene ein geringes Bewegungshemmnis.
\item Verwundbarkeit: Bist du gefesselt oder eingeschlossen, beginnst du in dieser Szene mit -10 auf die erste Fluchtprobe.
\end{itemize}
\noindent\textbf{Duigosz}\par
\begin{itemize}[leftmargin=*,nosep]
\item Vorteil: Erhalte +10 auf Proben für Handwerk, Reparatur oder den Betrieb von Maschinen.
\item Verwundbarkeit: Improvisierst du ohne geeignetes Werkzeug, sinkt der Effekt um eine Stufe (stark zu normal zu begrenzt).
\end{itemize}
\noindent\textbf{Katora}\par
\begin{itemize}[leftmargin=*,nosep]
\item Vorteil: In Dunkelheit erhältst du +10 auf Wahrnehmungsproben und Nahkampfproben.
\item Verwundbarkeit: In hartem direktem Licht erhältst du in jeder Szene -10 auf die erste Präzisionsprobe.
\end{itemize}
\noindent\textbf{Kroto'Chim}\par
\begin{itemize}[leftmargin=*,nosep]
\item Vorteil: Erhalte +10 auf Menschenkenntnis (Instinkt + Einfluss) und Proben zum Umgang mit Tieren.
\item Verwundbarkeit: In stark urbanen oder künstlichen Umgebungen erhältst du -10 auf Orientierung nach natürlichen Zeichen.
\end{itemize}
\noindent\textbf{Fraxut}\par
\begin{itemize}[leftmargin=*,nosep]
\item Vorteil: Bei schwachem Licht oder in Dunkelheit erhältst du +10 auf Wahrnehmungs- und Verfolgungsproben.
\item Verwundbarkeit: In hellem offenem Tageslicht erhältst du -10 auf Heimlichkeit, sofern du nicht zusätzliche Zeit für Deckungsvorbereitung aufwendest.
\end{itemize}
\noindent\textbf{Atiarel}\par
\begin{itemize}[leftmargin=*,nosep]
\item Vorteil: Erhalte +10 auf arkane Analyse (Geist + Magiekunde) und das Lesen von Ritualen.
\item Verwundbarkeit: Bei einer katastrophal misslungenen Magieprobe markierst du zusätzlich zum üblichen Rückschlag +1 Bürde.
\end{itemize}
\noindent\textbf{Silkanda}\par
\begin{itemize}[leftmargin=*,nosep]
\item Vorteil: Erhalte +10 auf gesellschaftliche Eleganz (Erscheinung + Einfluss) und diplomatische Erstkontakte.
\item Verwundbarkeit: Wirst du persönlich beleidigt oder entehrt, erhältst du -10 auf Proben zum Lösen oder Rückzug, bis du reagierst.
\end{itemize}
\noindent\textbf{Xordai}\par
\begin{itemize}[leftmargin=*,nosep]
\item Vorteil: Erhalte +10 auf Bergbau, Steinmetzarbeit und Orientierung unter Tage.
\item Verwundbarkeit: In tiefer Wildnis ohne Straßen oder Steinmarken erhältst du -10 auf Fernorientierung.
\end{itemize}
\noindent\textbf{Ancatir}\par
\begin{itemize}[leftmargin=*,nosep]
\item Vorteil: Erhalte +10 auf Heimlichkeit und Proben zur Vorbereitung von Fernangriffen aus der Deckung.
\item Verwundbarkeit: Wirst du ohne Vorbereitung in direkten Nahkampf gezwungen, erhält dein erster Nahkampfangriff in jeder Szene -10.
\end{itemize}
\noindent\textbf{Al Bah JiRa}\par
\begin{itemize}[leftmargin=*,nosep]
\item Vorteil: Erhalte +10 auf Anrufungen des Glaubens und formale Ritualführung.
\item Verwundbarkeit: Wenn du einen geschworenen Ritus oder ein Tabu brichst, kostet dich das sofort 1 Omen (oder +1 Bürde bei 0 Omen).
\end{itemize}
\noindent\textbf{Yavon}\par
\begin{itemize}[leftmargin=*,nosep]
\item Vorteil: Erhalte +10 auf Befehls-, Verhandlungs- und gesellschaftliche Proben im Zusammenhang mit Herrschaft.
\item Verwundbarkeit: Wird deine Autorität öffentlich untergraben, erhältst du in diesem Konflikt -10 auf die erste Reaktionsprobe.
\end{itemize}
\noindent\textbf{Nur'Tuk}\par
\begin{itemize}[leftmargin=*,nosep]
\item Vorteil: Erhalte +10 auf Standhaftigkeitsproben gegen Einschüchterung und Panikdruck.
\item Verwundbarkeit: In Situationen, die Geduld und Feinheit verlangen, erhältst du -10 auf die erste Täuschungs- oder Intrigenprobe.
\end{itemize}
\noindent\textbf{Hadewald}\par
\begin{itemize}[leftmargin=*,nosep]
\item Vorteil: Erhalte +10 auf Reitreisen, Tierführung und improvisierte Feldreparaturen.
\item Verwundbarkeit: In höfischen Etiketteszenen erhältst du -10 auf strenge Protokollproben, sofern dir kein anderer SC hilft.
\end{itemize}
\noindent\textbf{Meridian}\par
\begin{itemize}[leftmargin=*,nosep]
\item Vorteil: Erhalte +10 auf Proben zu Politik, Recht und historischen Präzedenzfällen.
\item Verwundbarkeit: Unter plötzlichem körperlichem Druck (Überfall, Verfolgung) erhältst du -10 auf die erste schnelle Handlungsprobe.
\end{itemize}
\noindent\textbf{Yadosia}\par
\begin{itemize}[leftmargin=*,nosep]
\item Vorteil: Erhalte +10 auf Heraldik, Auftritte und das Signalisieren sozialen Rangs.
\item Verwundbarkeit: In Unternehmungen, die auf Anonymität beruhen, sinkt dein Effekt um eine Stufe, sofern du keine vorbereitete Deckidentität annehmen kannst.
\end{itemize}
\noindent\textbf{Quitaron}\par
\begin{itemize}[leftmargin=*,nosep]
\item Vorteil: Erhalte +10 auf Überleben, Reiten und Ausdauerproben in der Wildnis.
\item Verwundbarkeit: In dichten Stadtvierteln erhältst du -10 auf Orientierung, sofern dich kein ortskundiger Führer begleitet.
\end{itemize}
\noindent\textbf{Gas'Danir}\par
\begin{itemize}[leftmargin=*,nosep]
\item Vorteil: Erhalte +10 auf Heimlichkeit, Infiltration und maritime Bewegung auf kurze Distanz.
\item Verwundbarkeit: In langen offenen Schlagabtauschen erhältst du -10 auf den zweiten und jeden weiteren Nahkampfaustausch derselben Szene.
\end{itemize}
\noindent\textbf{Gasdaria}\par
\begin{itemize}[leftmargin=*,nosep]
\item Vorteil: Erhalte +10 auf Verwaltung, Logistik und disziplinierte Koordination auf Distanz.
\item Verwundbarkeit: Wenn Pläne unerwartet zusammenbrechen, erhältst du -10 auf die erste Improvisationsprobe.
\end{itemize}
\noindent\textbf{Toran}\par
\begin{itemize}[leftmargin=*,nosep]
\item Vorteil: Einmal pro Szene erhält eine soziale Einflussprobe +10, wenn sie Schulden, Rang oder Zwangspolitik ausnutzt.
\item Verwundbarkeit: Verweigerst du eine direkte Herausforderung durch einen Rivalen, kostet dich das 1 Omen (oder +1 Bürde bei 0 Omen).
\end{itemize}


\section{Wege}\label{sec:wege}
\noindent\textbf{Spion}\par
\begin{itemize}[leftmargin=*,nosep]
\item Vorteil: Erhalte +10 auf Infiltration, Überwachung und verdeckte Kommunikationsproben.
\item Verwundbarkeit: Ein bekannter Rivale kann dir -10 auf die erste soziale Probe einer Szene geben, in der deine Identität offengelegt ist.
\end{itemize}
\noindent\textbf{Soldat}\par
\begin{itemize}[leftmargin=*,nosep]
\item Vorteil: Erhalte +10 auf koordinierte Kampfmanöver und Proben auf Schlachtfelddisziplin.
\item Verwundbarkeit: Bei direkten widersprüchlichen Befehlen oder Befehlszusammenbruch markierst du +1 Bürde oder erhältst in dieser Runde -10 auf Proben, die über deine Handlungsreihenfolge entscheiden.
\end{itemize}
\noindent\textbf{Geistliche}\par
\begin{itemize}[leftmargin=*,nosep]
\item Vorteil: Erhalte +10 auf Riten, Liturgie und Überzeugungsproben aus der Lehre heraus.
\item Verwundbarkeit: Handelst du offen gegen deine erklärte Lehre, verlierst du 1 Omen (oder markierst +1 Bürde bei 0 Omen).
\end{itemize}
\noindent\textbf{Waldläufer}\par
\begin{itemize}[leftmargin=*,nosep]
\item Vorteil: Erhalte +10 auf Verfolgung, Feldhandwerk und die Vorbereitung von Fernhaltungen im wilden Gelände.
\item Verwundbarkeit: In dichten Stadtvierteln erhältst du -10 auf Orientierung, sofern dir kein ortskundiger Führer hilft.
\end{itemize}
\noindent\textbf{Schmied}\par
\begin{itemize}[leftmargin=*,nosep]
\item Vorteil: Erhalte +10 auf Schmieden, Reparatur und Proben zur Beurteilung von Metall.
\item Verwundbarkeit: Ohne geeignetes Werkzeug oder Schmiede sinkt der Effekt bei Ausrüstungsreparaturen um eine Stufe.
\end{itemize}
\noindent\textbf{Boxer}\par
\begin{itemize}[leftmargin=*,nosep]
\item Vorteil: Erhalte +10 auf unbewaffneten Nahkampf und Proben auf Schmerzausdauer.
\item Verwundbarkeit: Gegen Reichweitenwaffen auf offenem Boden erhältst du in jeder Szene -10 auf den ersten Nahkampfaustausch.
\end{itemize}
\noindent\textbf{Schreiber}\par
\begin{itemize}[leftmargin=*,nosep]
\item Vorteil: Erhalte +10 auf Aufzeichnungen, Schriften und Proben zur Echtheit von Dokumenten.
\item Verwundbarkeit: Unter unmittelbarem Zeitdruck erhältst du in der Szene -10 auf die erste Schreib- oder Analyseprobe.
\end{itemize}
\noindent\textbf{Söldner}\par
\begin{itemize}[leftmargin=*,nosep]
\item Vorteil: Erhalte +10 bei Vertragsverhandlungen, Drohgebärden und praktischer Scharmützeltaktik.
\item Verwundbarkeit: Werden Bezahlung, Bedingungen oder Loyalität öffentlich infrage gestellt, erhältst du -10 auf die erste Standhaftigkeitsprobe auf Moral oder Mut.
\end{itemize}
\noindent\textbf{Seefahrer}\par
\begin{itemize}[leftmargin=*,nosep]
\item Vorteil: Erhalte +10 auf Seemannschaft, Tauwerk und Navigationsproben im Sturm.
\item Verwundbarkeit: Bei Überlandreisen durch Berge oder Wüste erhältst du an jedem Tag -10 auf die erste Orientierungsprobe.
\end{itemize}
\noindent\textbf{Gelehrte}\par
\begin{itemize}[leftmargin=*,nosep]
\item Vorteil: Erhalte +10 auf Geschichte, Mythen und Proben zur Zusammenführung von Erkenntnissen.
\item Verwundbarkeit: Bei plötzlicher Gewalt erhältst du in dem Konflikt -10 auf die erste Reaktionsprobe.
\end{itemize}
\noindent\textbf{Abenteurer}\par
\begin{itemize}[leftmargin=*,nosep]
\item Vorteil: Einmal pro Szene ignorierst du ein geringes Hindernis der Fortbewegung und erhältst +10 auf Proben zu Erkundungsrisiken.
\item Verwundbarkeit: Bist du überdehnt (wenig Vorräte oder keine Rast), markierst du nach einer misslungenen Gefahrenprobe +1 Bürde.
\end{itemize}
\noindent\textbf{Barde}\par
\begin{itemize}[leftmargin=*,nosep]
\item Vorteil: Erhalte +10 auf Auftritte, soziale Einflussnahme und Szenen, in denen du Moral stärkst.
\item Verwundbarkeit: Wirst du öffentlich verspottet oder diskreditiert, erhältst du danach -10 auf die erste Überzeugungsprobe.
\end{itemize}
\noindent\textbf{Bote}\par
\begin{itemize}[leftmargin=*,nosep]
\item Vorteil: Erhalte +10 auf schnelle Reisen und Proben auf Eilzustellung unter Druck.
\item Verwundbarkeit: Bist du mit schwerer Ausrüstung beladen, erhältst du -10 auf Verfolgungs- oder Fluchtbewegung.
\end{itemize}
\noindent\textbf{Händler}\par
\begin{itemize}[leftmargin=*,nosep]
\item Vorteil: Erhalte +10 auf Feilschen, Wertschätzung und Proben auf Handelsnetzwerke.
\item Verwundbarkeit: In Geschenkpflicht- oder Tabuökonomien erhältst du -10 auf die erste Handelsverhandlungsprobe.
\end{itemize}
\noindent\textbf{Dieb}\par
\begin{itemize}[leftmargin=*,nosep]
\item Vorteil: Erhalte +10 auf heimliches Eindringen, Schlossöffnung und Proben auf Fingerfertigkeit.
\item Verwundbarkeit: Nach einer misslungenen Heimlichkeitsprobe erhältst du in derselben Szene -10 auf deine nächste soziale Tarnprobe.
\end{itemize}
\noindent\textbf{Bestatter}\par
\begin{itemize}[leftmargin=*,nosep]
\item Vorteil: Erhalte +10 auf Totenriten, Leichenumgang und ruhiges Handeln in Szenen des Todes.
\item Verwundbarkeit: Lebende Menschen misstrauen deinem Ton; in festlichen Situationen erhältst du -10 auf die erste charmante soziale Probe.
\end{itemize}
\noindent\textbf{Meuchler}\par
\begin{itemize}[leftmargin=*,nosep]
\item Vorteil: Erhalte +10 auf geplante Tötungsversuche und Proben zur Vorbereitung durch Zielstudium.
\item Verwundbarkeit: In langem offenem Kampf erhältst du ab der zweiten Runde -10 auf Präzisionsangriffe.
\end{itemize}
\noindent\textbf{Medicus}\par
\begin{itemize}[leftmargin=*,nosep]
\item Vorteil: Erhalte +10 auf Diagnose, Chirurgie und Notbehandlungsproben.
\item Verwundbarkeit: Musst du ohne saubere Werkzeuge arbeiten, fügt eine misslungene Heilprobe eine zusätzliche Komplikation hinzu.
\end{itemize}
\noindent\textbf{Ritter}\par
\begin{itemize}[leftmargin=*,nosep]
\item Vorteil: Erhalte +10 auf berittene Kampfkommandos, Ehrenduelle und formale Autoritätshaltung.
\item Verwundbarkeit: Das Ignorieren einer erklärten Herausforderung kostet 1 Omen (oder +1 Bürde bei 0 Omen).
\end{itemize}
\noindent\textbf{Reiter}\par
\begin{itemize}[leftmargin=*,nosep]
\item Vorteil: Erhalte +10 auf Reiten, berittene Verfolgung und Proben zur Tierführung.
\item Verwundbarkeit: In engen Innenräumen oder Tunneln erhältst du in jeder Szene -10 auf die erste Bewegungsprobe.
\end{itemize}
\noindent\textbf{Medium}\par
\begin{itemize}[leftmargin=*,nosep]
\item Vorteil: Erhalte +10, wenn du Omen, Echos oder Anzeichen geisterhafter Präsenz deutest.
\item Verwundbarkeit: Bei katastrophal misslungenem okkultem Kontakt markierst du zusätzlich zum normalen Nachhall +1 Bürde.
\end{itemize}
\noindent\textbf{Jäger}\par
\begin{itemize}[leftmargin=*,nosep]
\item Vorteil: Erhalte +10 auf Verfolgung, die Vorbereitung von Schüssen und das Lesen von Beute in Wildnisszenen.
\item Verwundbarkeit: In höfischen oder juristischen Verhandlungsszenen erhältst du -10 auf die erste Einflussprobe.
\end{itemize}
\noindent\textbf{Bader}\par
\begin{itemize}[leftmargin=*,nosep]
\item Vorteil: Erhalte +10 auf grobe Traumaversorgung und Proben zur Stabilisierung im Feld.
\item Verwundbarkeit: Bei feiner langwieriger Behandlung erhältst du -10 auf die erste präzise medizinische Probe.
\end{itemize}
\noindent\textbf{Abdecker}\par
\begin{itemize}[leftmargin=*,nosep]
\item Vorteil: Erhalte +10 auf Einschüchterung, die Verarbeitung von Kadavern und Ausdauer bei grimmiger Arbeit.
\item Verwundbarkeit: In etikettelastigen sozialen Szenen erhältst du -10 auf die erste Diplomatieprobe.
\end{itemize}
\noindent\textbf{Wirt}\par
\begin{itemize}[leftmargin=*,nosep]
\item Vorteil: Erhalte +10 auf Gerüchtejagd, das Lesen von Menschenmengen und Proben zur Gastfreundschaftslogistik.
\item Verwundbarkeit: Beginnt eine Szene mit offener Gewalt, erhältst du -10 auf die erste Befehls- oder Deeskalationsprobe.
\end{itemize}


\section{Bindungen}\label{sec:bindungen}
\noindent\textbf{Verfallene Kapelle}\par
\begin{itemize}[leftmargin=*,nosep]
\item Vorteil: Erhalte +10 auf Proben mit Ritus, Mythen oder Standhaftigkeit, die mit geweihten Ruinen, Eiden oder Begräbnisriten zusammenhängen.
\item Verwundbarkeit: Wenn du in einem heiligen Ort die Hilflosen zurücklässt, markierst du +1 Bürde und die Bindung ist erschöpft.
\end{itemize}
\noindent\textbf{Schwarzsalz-Kompanie}\par
\begin{itemize}[leftmargin=*,nosep]
\item Vorteil: Erhalte +10 auf Vertragshebel, Söldnerverhandlungen oder Schlachtfeldkoordination mit Veteranen.
\item Verwundbarkeit: Verweigerst du einem früheren Mitglied der Kompanie unter Druck Hilfe, verlierst du 1 Omen (oder markierst +1 Bürde bei 0 Omen).
\end{itemize}
\noindent\textbf{Trauerschwester}\par
\begin{itemize}[leftmargin=*,nosep]
\item Vorteil: Einmal pro Szene kannst du eine soziale Konsequenz oder Druckkonsequenz um eine Stufe senken, wenn du handelst, um Verwandte zu schützen.
\item Verwundbarkeit: Wird deine Schwester bedroht, erhält deine erste Handlung ohne Unterstützungsabsicht in dieser Szene -10.
\end{itemize}
\noindent\textbf{Alter Mentor}\par
\begin{itemize}[leftmargin=*,nosep]
\item Vorteil: Erhalte einmal pro Szene +10 auf eine Vorbereitungs-, Recherche- oder taktische Aufbauprobe, nachdem du die Lehre deines Mentors benannt hast.
\item Verwundbarkeit: Ignorierst du im Spiel eine bekannte Warnung deines Mentors, erhält deine nächste passende Probe -10.
\end{itemize}
\noindent\textbf{Stadtunterkrypta}\par
\begin{itemize}[leftmargin=*,nosep]
\item Vorteil: Erhalte +10 auf Heimlichkeit, Orientierung und okkultes Erinnern in Katakomben, Abwasserläufen und Grabgängen.
\item Verwundbarkeit: In offenen Tageslichtvierteln erhältst du in jeder Szene -10 auf die erste Wahrnehmungs- oder Orientierungsprobe.
\end{itemize}


\section{Male}\label{sec:male}
\noindent\textbf{Verbrannte Hände}\par
\begin{itemize}[leftmargin=*,nosep]
\item Vorteil: Erhalte +10 auf Proben zu Schmerzausdauer und Hitzegefahren.
\item Verwundbarkeit: Bei feinen Handarbeiten erhältst du -10, sofern du nicht geeignetes Werkzeug oder zusätzliche Zeit einsetzt.
\end{itemize}
\noindent\textbf{Eidmal}\par
\begin{itemize}[leftmargin=*,nosep]
\item Vorteil: Erhalte +10 auf Befehl oder Einschüchterung, wenn du dein geschworenes Mal anrufst.
\item Verwundbarkeit: Wenn du dein gesprochenes Versprechen in der Szene brichst, kostet dich das 1 Omen (oder +1 Bürde bei 0 Omen).
\end{itemize}
\noindent\textbf{Fehlendes Auge}\par
\begin{itemize}[leftmargin=*,nosep]
\item Vorteil: Erhalte +10 auf das Lesen von Hinterhaltswinkeln und das Vorausahnen von Bedrohungen.
\item Verwundbarkeit: Bei Präzisionsproben auf lange Distanz erhältst du -10.
\end{itemize}
\noindent\textbf{Gebrochener Kiefer}\par
\begin{itemize}[leftmargin=*,nosep]
\item Vorteil: Erhalte +10 auf das Widerstehen von Zwang, Panik oder Verhördruck.
\item Verwundbarkeit: In einer formalen Szene erhältst du -10 auf die erste Charme- oder Diplomatieprobe.
\end{itemize}
\noindent\textbf{Nachtzittern}\par
\begin{itemize}[leftmargin=*,nosep]
\item Vorteil: Erhalte +10, wenn du im Lager oder in der Ruhe Gefahr erspürst.
\item Verwundbarkeit: Nach einer Nacht ohne Sicherheit beginnst du die nächste Szene mit +1 Bürde.
\end{itemize}
\noindent\textbf{Pockennarbige Haut}\par
\begin{itemize}[leftmargin=*,nosep]
\item Vorteil: Erhalte +10 auf Einschüchterung gegen Abergläubische.
\item Verwundbarkeit: Bei Etikette oder höfischen Ersteindrücken erhältst du -10.
\end{itemize}
\noindent\textbf{Hinkendes Bein}\par
\begin{itemize}[leftmargin=*,nosep]
\item Vorteil: Erhalte +10 auf Standfestigkeit, wenn du unter Druck Boden hältst.
\item Verwundbarkeit: Bei Verfolgungen oder schnellem Umpositionieren erhältst du -10 auf die erste Bewegungsprobe.
\end{itemize}
\noindent\textbf{Zerschlagene Knöchel}\par
\begin{itemize}[leftmargin=*,nosep]
\item Vorteil: Erhalte +10 auf unbewaffnete Schläge und Schmerzwiderstand.
\item Verwundbarkeit: Bei feiner Handwerks- oder Heilarbeit sinkt der Effekt um eine Stufe.
\end{itemize}
\noindent\textbf{Rituelle Narbung}\par
\begin{itemize}[leftmargin=*,nosep]
\item Vorteil: Erhalte +10 auf okkulte Wiedererkennung und das Erinnern von Riten.
\item Verwundbarkeit: Unheilvolle Zeichen und Gesänge können dich ablenken: Die erste Konzentrationsprobe erhält -10.
\end{itemize}
\noindent\textbf{Würgemal am Hals}\par
\begin{itemize}[leftmargin=*,nosep]
\item Vorteil: Erhalte +10 auf Entfesselungskunst und Proben zum Brechen von Fesseln.
\item Verwundbarkeit: In Szenen mit erstickendem Rauch, Würgegas oder tiefem Wasser erhältst du -10 auf die erste Ausdauerprobe.
\end{itemize}
\noindent\textbf{Versengte Lungen}\par
\begin{itemize}[leftmargin=*,nosep]
\item Vorteil: Erhalte +10 auf das Widerstehen von Rauch- und Ascheeinwirkung.
\item Verwundbarkeit: Sprinten oder lange Anstrengung geben in einer Szene -10 auf die zweite und jede weitere körperliche Probe.
\end{itemize}
\noindent\textbf{Verstümmeltes Ohr}\par
\begin{itemize}[leftmargin=*,nosep]
\item Vorteil: Erhalte +10 auf Lippenlesen und das Deuten von Körpersprache in Menschenmengen.
\item Verwundbarkeit: In lärmendem Chaos mit schlechter Sicht erhältst du -10 auf hörbasierte Wahrnehmung.
\end{itemize}
\noindent\textbf{Angebrochene Rippen}\par
\begin{itemize}[leftmargin=*,nosep]
\item Vorteil: Erhalte +10 darauf, stumpfe Gewalt zu überstehen.
\item Verwundbarkeit: Beim Klettern, Ringen oder schweren Heben erhältst du -10, sofern du dich nicht abstützt.
\end{itemize}
\noindent\textbf{Eisenstift in der Schulter}\par
\begin{itemize}[leftmargin=*,nosep]
\item Vorteil: Erhalte +10 auf Waffenhalt und das Widerstehen gegen Entwaffnung.
\item Verwundbarkeit: Schläge oder Würfe über Kopf erhalten ohne Vorbereitung -10.
\end{itemize}
\noindent\textbf{Erfrorene Finger}\par
\begin{itemize}[leftmargin=*,nosep]
\item Vorteil: Erhalte +10 auf Überlebensproben in extremer Kälte.
\item Verwundbarkeit: In nasskalten Bedingungen erhältst du -10 auf die erste präzise Handprobe.
\end{itemize}
\noindent\textbf{Ausgebrannte Wunde}\par
\begin{itemize}[leftmargin=*,nosep]
\item Vorteil: Einmal pro Szene ignorierst du eine geringe Blutungskonsequenz.
\item Verwundbarkeit: Magische Heilung an dir ist instabil: Eine misslungene Heilprobe fügt eine zusätzliche Komplikation hinzu.
\end{itemize}
\noindent\textbf{Vom Peitschenhieb gezeichneter Rücken}\par
\begin{itemize}[leftmargin=*,nosep]
\item Vorteil: Erhalte +10 auf Gewaltmärsche und das Tragen von Last.
\item Verwundbarkeit: Plötzliche Treffer auf den Oberkörper geben -10 auf unmittelbare Reaktionsproben.
\end{itemize}
\noindent\textbf{Salz in den Adern}\par
\begin{itemize}[leftmargin=*,nosep]
\item Vorteil: Erhalte +10 auf das Widerstehen von Gift und Krankheit.
\item Verwundbarkeit: Starke alchemische oder medizinische Geschmäcker können Würgereiz auslösen: Die erste Probe auf eingenommene Mittel erhält -10.
\end{itemize}
\noindent\textbf{Überlebende des Grabfiebers}\par
\begin{itemize}[leftmargin=*,nosep]
\item Vorteil: Erhalte +10 auf Standhaftigkeitsproben gegen Untote und Leichenbedrohungen.
\item Verwundbarkeit: In der Nähe offener Gräber oder Beinhaufen markierst du +1 Bürde, sofern dir nicht Körper + Standhaftigkeit gelingt.
\end{itemize}
\noindent\textbf{Gespaltene Zungennarbe}\par
\begin{itemize}[leftmargin=*,nosep]
\item Vorteil: Erhalte +10 auf Codesprache, Gaunersprache und rituelle Formulierungen.
\item Verwundbarkeit: In formaler Rhetorik und öffentlicher Rede erhältst du -10 auf die erste Überzeugungsprobe.
\end{itemize}
